\documentclass[11pt]{article}

\usepackage[letterpaper,margin=1in]{geometry}
\usepackage{amsmath,amssymb,amsthm,mathtools}
\usepackage{array,booktabs,enumitem,microtype,placeins}
\usepackage[hidelinks]{hyperref}

\theoremstyle{definition}
\newtheorem{definition}{Definition}[section]
\newtheorem{example}[definition]{Example}
\theoremstyle{plain}
\newtheorem{theorem}[definition]{Theorem}
\newtheorem{proposition}[definition]{Proposition}
\newtheorem{corollary}[definition]{Corollary}
\newtheorem{lemma}[definition]{Lemma}
\theoremstyle{remark}
\newtheorem{remark}[definition]{Remark}

\newcommand{\True}{\mathrm{T}}
\newcommand{\False}{\mathrm{F}}
\newcommand{\Top}{\operatorname{Top}}
\newcommand{\Int}{\operatorname{Int}}
\newcommand{\Cl}{\operatorname{Cl}}
\newcommand{\bd}{\partial}
\newcommand{\Pow}{\mathcal P}
\newcommand{\restr}{\mathord{|}}
\newcommand{\ind}{\mathbf 1}
\newcommand{\modelsF}{\models}

\title{Propositional Logical Topology:\\Observation, Distinguishability, and Inference}
\author{Brian Theory}
\date{Core manuscript v0.8.1, September 2026}

\hypersetup{
  pdftitle={Propositional Logical Topology: Observation, Distinguishability, and Inference},
  pdfauthor={Brian Theory}
}

\begin{document}
\maketitle

\begin{abstract}
Fix a propositional language over finitely many variables and regard its truth
assignments as points.  A selected family of formulas acts as a family of
positive-certification tests.  The topology generated by their truth regions
records which logical states can be distinguished by those tests.  We show that
topological indistinguishability is equality of observation signatures, that
the specialization preorder is inclusion of positive observations, and that a
directed loss quasi-pseudometric generates the same topology.  Its sum with the
reverse distance is weighted Hamming distance on observation signatures.  We
then separate two forms of information acquisition: refinement of observable
tests and restriction of possible valuations by premises.  Classical inference
remains inclusion of the common premise region in the conclusion region; the
observation topology organizes certification, subspaces, and continuous maps
around that inclusion.  The contribution is an expository synthesis of standard
finite Boolean and Alexandrov-space constructions, with emphasis on distinctions
that are easily conflated in logical-topological notation.
\end{abstract}

\tableofcontents

\section{Introduction}

Propositional formulas can be represented by the truth assignments on which
they hold.  If every such subset of a finite valuation space is declared open,
the resulting topology is discrete and contains no structure beyond classical
truth-set semantics.  A nontrivial topological treatment begins instead with a
restricted family of formulas that can be used as observations.  Two truth
assignments may then be logically different but observationally
indistinguishable.

This paper develops five linked ideas:
\begin{enumerate}[label=(\roman*)]
\item complete truth assignments give the atomic partition for finite Boolean
      semantics;
\item selected observable formulas determine distinguishability and
      indistinguishability;
\item the specialization preorder records asymmetric logical separation;
\item XOR of observation signatures gives symmetric disagreement, while a
      source mask gives directed loss; and
\item adding observations refines the topology, while adding premises
      restricts the current space of possibilities.
\end{enumerate}

The intended reader knows propositional truth tables and basic set notation but
may be meeting finite topology, specialization order, or topological semantics
for the first time.  The expository aim is to use one finite observation table
to show which constructions share data and which do not: in particular,
truth-region intersection versus intersection of topologies, complete semantic
profiles versus partial certificates, and an indexed observation presentation
versus the topology it generates.

The qualitative construction can be summarized as
\[
  \Phi\longmapsto o_\Phi\longmapsto\tau_\Phi
  \longleftrightarrow\preceq_\Phi.
\]
Given positive weights $\mathbf a=(a_1,\ldots,a_m)$, the same presentation also
gives
\[
  (\Phi,\mathbf a)\longmapsto d^+_{\Phi,\mathbf a}
  \longmapsto (\preceq_\Phi,\tau_\Phi).
\]
Here $\Phi$ is an indexed family of observable formulas, $o_\Phi$ is the observation
signature, $\tau_\Phi$ is the generated topology, $\preceq_\Phi$ is its
specialization preorder, and $d_\Phi^+$ is a directed logical distance.  The
last arrow is not reversible: the topology determines its specialization
preorder on this finite carrier, but it does not determine the chosen weights
or observation presentation.

The Boolean representation underlying the truth-region map is classical and
is closely related to the representation of Boolean algebras by fields of sets
\cite{Stone1936}.  Topological semantics for logical operations has a long
history, notably in the work of McKinsey and Tarski \cite{McKinseyTarski1944}.
Finite topologies and their order-theoretic structure are classical as well
\cite[pp.~325--327]{Stong1966}; see also \cite{Alexandroff1937}.
Sierpi\'nski-valued predicates provide a standard description of opens and
positive observation \cite[Ch.~1]{Vickers1989}; see also
\cite[Example~1.12(3)]{Vickers2007}.
Observation frames provide a related abstract separation between observable
predicates and all predicates \cite{BonsangueJacobsKok1995}.  The purpose here
is to assemble these ideas around partitioned propositional states, asymmetric
distinguishability, signature distance, refinement, and inference.  No new
duality or completeness theorem is claimed.

Two further precedents use essentially the same object--attribute data.  In
Pawlak's information systems, equality of attribute descriptions induces an
indiscernibility partition, whose lower and upper approximations are interior
and closure in the associated partition topology
\cite[Sec.~2.2, pp.~343--344]{Pawlak1982}.  That
construction agrees here with two-sided observation; positive observation may
instead retain strict specialization comparisons.  In formal concept analysis,
the incidence $v\mathrel I i$ defined by $v\models\varphi_i$ is a formal context,
and a finite conjunction of attributes has the corresponding extent
\cite[Sec.~1.1]{GanterWille1999}.  Our topology also admits arbitrary unions of such
extents, which need not themselves be extents, so its open-set lattice is not in
general the concept lattice.  For $\Phi=(X,Y)$, for instance, the three-element
set of valuations $\{(1,0),(1,1),(0,1)\}$ is open but is not an extent of the
context with just the attributes $X$ and $Y$.

\section{Valuations, truth regions, and partitions}

Let $P=\{p_1,\ldots,p_n\}$ be a finite set of propositional variables and let
\[
  \Omega_P=\{0,1\}^{P}
\]
be the set of valuations.  For $v\in\Omega_P$, the expression
$\mathcal V_v(\mathcal L)$ is the value of a formula $\mathcal L$ at $v$, under
the usual classical truth tables.  We identify $\True$ with $1$ and $\False$
with $0$.

\subsection{Formula tuples, truth vectors, and truth regions}

\begin{definition}[Formula tuple and truth vector]
A \emph{formula tuple} of arity $m$ is an ordered tuple of formulas
\[
 \mathbf L=(\mathcal L_1,\ldots,\mathcal L_m).
\]
At a valuation $v\in\Omega_P$, its \emph{truth vector} is the Boolean vector
\[
 \mathcal V_v(\mathbf L)
 :=\bigl(\mathcal V_v(\mathcal L_1),\ldots,
          \mathcal V_v(\mathcal L_m)\bigr)
 \in\{0,1\}^m.
\]
\end{definition}

A formula tuple is therefore syntactic, whereas a truth vector is one of its
possible semantic values.  Neither should be confused with a valuation
$v\in\Omega_P$, which assigns truth values to every variable in the fixed
language and serves as a point of the topological carrier.

\begin{definition}[Truth set or truth region]
For a formula $\mathcal L$, define
\[
  \varsigma(\mathcal L)
  :=\{v\in\Omega_P:\mathcal V_v(\mathcal L)=\True\}.
\]
The notation $\mathcal V_v^{\True}(\mathcal L)$ abbreviates
$v\in\varsigma(\mathcal L)$, and
$\mathcal V_v^{\False}(\mathcal L)$ abbreviates
$v\notin\varsigma(\mathcal L)$.
\end{definition}

The terms \emph{truth set} and \emph{truth region} refer to this same subset of
$\Omega_P$.  We use ``region'' when emphasizing its role inside a topological
space.  The superscripts $\True$ and $\False$ above denote truth assertions;
they do not define additional numerical evaluation maps.

Using a fixed carrier ensures that $\varsigma$ depends only on the logical
equivalence class of $\mathcal L$.  It obeys
\begin{align*}
 \varsigma(\mathcal L_1\wedge\mathcal L_2)
   &=\varsigma(\mathcal L_1)\cap\varsigma(\mathcal L_2),\\
 \varsigma(\mathcal L_1\vee\mathcal L_2)
   &=\varsigma(\mathcal L_1)\cup\varsigma(\mathcal L_2),\\
 \varsigma(\neg\mathcal L)
   &=\Omega_P\setminus\varsigma(\mathcal L).
\end{align*}

Write $\mathcal L_1\Updownarrow\mathcal L_2$ for Boolean exclusive OR.  For two
formulas, their union separates into an overlap and an exclusive remainder:
\[
 \varsigma(\mathcal L_1\vee\mathcal L_2)
 =\varsigma(\mathcal L_1\wedge\mathcal L_2)
  \mathbin{\dot\cup}
  \varsigma(\mathcal L_1\Updownarrow\mathcal L_2).
\]
Equivalently,
\[
 \mathcal L_1\vee\mathcal L_2
 \equiv
 (\mathcal L_1\wedge\mathcal L_2)
 \Updownarrow
 (\mathcal L_1\Updownarrow\mathcal L_2),
\]
because the overlap and exclusive remainder are disjoint.  The complete
two-formula truth-pattern partition is
\begin{align*}
 &\mathcal L_1\wedge\mathcal L_2,
 &&\mathcal L_1\wedge\neg\mathcal L_2,\\
 &\neg\mathcal L_1\wedge\mathcal L_2,
 &&\neg\mathcal L_1\wedge\neg\mathcal L_2,
\end{align*}
after empty cells are removed.  Thus XOR identifies the two exclusive cells;
it is not a conversion of the conjunction into a disjunction.

For every $v\in\Omega_P$, define its complete minterm
\[
  m_v:=\bigwedge_{i=1}^n \ell_i,
  \qquad
  \ell_i=
  \begin{cases}
    p_i,&v(p_i)=1,\\
    \neg p_i,&v(p_i)=0.
  \end{cases}
\]
The family $\{\varsigma(m_v):v\in\Omega_P\}$ is the singleton partition of
$\Omega_P$.  Every formula has the canonical complete disjunctive form
\[
  \mathcal L\equiv
  \bigvee_{v\in\varsigma(\mathcal L)}m_v,
\]
where the empty disjunction is $\False$ and the empty conjunction is $\True$.
Thus the complete minterms give the atoms of the finite Boolean algebra
$\Pow(\Omega_P)$.  They form a topological basis only when the topology is
discrete.

For two variables use the cell labels
\[
\begin{array}{ll}
 A=X\wedge\neg Y, & B=X\wedge Y,\\
 C=\neg X\wedge Y, & D=\neg X\wedge\neg Y,
\end{array}
\qquad
\Omega_{X,Y}=\{A,B,C,D\}.
\]
Each letter denotes the unique valuation satisfying the displayed minterm; the
same letter may also label that singleton truth cell when the type is clear.
For example,
\[
 \varsigma(X)=\{A,B\},\quad
 \varsigma(Y)=\{B,C\},\quad
 \varsigma(X\Rightarrow Y)=\{B,C,D\}.
\]

\section{Observation-generated spaces}

The topology should record a limited observational vocabulary rather than all
definable subsets.

\begin{definition}[Observation family and signature]
Let $I:=\{1,\ldots,m\}$ and let the observations be the indexed finite family
\[
  \Phi=(\varphi_i)_{i\in I}.
\]
Put
\[
  b_i(v):=\ind[v\in\varsigma(\varphi_i)]
\]
and define the observation signature
\[
  o_\Phi(v):=(b_i(v))_{i\in I}\in\{0,1\}^I.
\]
When subset notation is used, we identify this vector with its support
\[
  o_\Phi(v)=\{i\in I:v\modelsF\varphi_i\}\subseteq I.
\]
Different indices may carry logically equivalent formulas, or even literal
repetitions, and may receive distinct weights.
\end{definition}

The truth vector of the indexed formula tuple $\Phi$ at $v$ is
exactly the observation signature:
\[
 \mathcal V_v(\Phi)=o_\Phi(v).
\]
The signature notation emphasizes that these coordinates are the observations
available to the topology.

The evidence convention is asymmetric: a subbasic event records positive
certification that $\varphi_i$ is true.  Failure to certify $\varphi_i$ is not
itself a positive certificate for $\neg\varphi_i$.  If both outcomes of every
Boolean test are operationally available, the appropriate family is
$\Phi^{\pm}$ from Proposition~\ref{prop:two-sided-observation} below, and the
resulting topology is the two-sided partition topology.  Here ``positive'' is
relative to the selected tests; a selected formula may itself contain a
negation.

The signature is a complete semantic profile: $b_i(v)=0$ asserts that
$\varphi_i$ is false at $v$.  It does not mean that a certificate has merely
failed to arrive.  The topology describes the affirmative distinctions allowed
by the selected interface; it does not by itself model time, partial transcripts,
or delivery of certificates.  Likewise, the distances in Section~6 are defined
on complete profiles and are not automatically measurable from an unfinished
positive-only transcript.

For a family $\mathcal G$ of subsets of $\Omega_P$, let
$\Top(\mathcal G)$ denote the smallest topology on $\Omega_P$ containing
$\mathcal G$.

\begin{definition}[Observation-generated topology and ProLT]
The topology generated by $\Phi$ is
\[
  \tau_\Phi:=\Top\{\varsigma(\varphi_i):i\in I\}
\]
on the fixed carrier $\Omega_P$.  The displayed truth regions are a subbasis;
their finite intersections form a basis.  The associated \emph{propositional
logical topological space}, abbreviated \emph{ProLT}, is the pair
$(\Omega_P,\tau_\Phi)$.
\end{definition}

We also call $(\Omega_P,\tau_\Phi)$ the \emph{observation space generated by
$\Phi$}.  Its points are valuations and its designated subbasic opens are truth
sets.  The carrier is always displayed because a family of subsets is not a
topology until its underlying set is fixed.  The indexed family $\Phi$ and any
assigned weights are presentation data; they are not part of the topology
$\tau_\Phi$ and need not be recoverable from it.

Because complete minterms define every subset of the finite carrier, every
topology on $\Omega_P$ can be presented by choosing formulas for its open sets.
Thus the propositional language supplies a finite presentation and
interpretation; observation-generated ProLTs are not a new restricted class of
finite topological spaces.

The empty intersection of subbasic opens is $\Omega_P$.  Hence a convenient
basis is
\[
  \mathcal B_\Phi
  =\left\{
    \bigcap_{i\in J}\varsigma(\varphi_i):J\subseteq I
   \right\}.
\]

Equivalently, give $\mathbb S=\{0,1\}$ the Sierpi\'nski topology
$\{\varnothing,\{1\},\mathbb S\}$.  Then $\tau_\Phi$ is the initial topology of
the coordinate maps $b_i:\Omega_P\to\mathbb S$, or of their product
$o_\Phi:\Omega_P\to\mathbb S^I$.  This is the standard topological formulation
of positive observable predicates \cite[Example~1.12(3)]{Vickers2007}.

Order the Boolean cube $\{0,1\}^I$ coordinatewise, or equivalently order its
subsets by inclusion.

\begin{theorem}[Signature representation]
The opens of $\tau_\Phi$ are precisely the inverse images of upward-closed sets of
realized observation signatures:
\[
  \tau_\Phi
  =\{o_\Phi^{-1}(U):U\subseteq o_\Phi(\Omega_P)
       \text{ is upward closed}\}.
\]
\end{theorem}

\begin{proof}
For a realized signature $s$, its principal upward set is
\[
  \mathord\uparrow s
  =\{t:s\subseteq t\}.
\]
Its inverse image is the basic open
\[
  o_\Phi^{-1}(\mathord\uparrow s)
  =\bigcap_{i\in s}\varsigma(\varphi_i).
\]
Every upward-closed set in a finite poset is the union of the principal upward
sets of its members, so its inverse image is open.  Conversely, a basic open
indexed by any $J\subseteq I$ is the inverse image of the upward set
\[
 \{s\in o_\Phi(\Omega_P):J\subseteq s\}.
\]
Unions of basic opens give the reverse inclusion.
\end{proof}

\subsection{Partitions induced by observations}

Define
\[
  v\sim_\Phi w
  \quad\Longleftrightarrow\quad
  o_\Phi(v)=o_\Phi(w).
\]
This equivalence relation partitions $\Omega_P$ into observational states.

\begin{definition}[Partition topology]
For a partition $\Pi$ of a set $S$, define
\[
  \tau_\Pi:=\left\{\bigcup\mathcal A:\mathcal A\subseteq\Pi\right\}.
\]
Its blocks are both open and closed.
\end{definition}

An arbitrary observation topology need not be a partition topology.  The
difference is essential: partition topologies describe symmetric
indistinguishability classes, whereas general finite topologies can also record
one-way observational dependence.

\begin{proposition}[Positive and two-sided observation]\label{prop:two-sided-observation}
Let $\Phi^{\pm}$ be indexed by $I\times\{+,-\}$, with
$\varphi_{(i,+)}=\varphi_i$ and $\varphi_{(i,-)}=\neg\varphi_i$.
Then $\tau_{\Phi^{\pm}}$ is the partition topology whose blocks are the
$\sim_\Phi$-classes.
\end{proposition}

\begin{proof}
For every realized signature $s$, its fiber is
\[
 o_\Phi^{-1}(s)
 =\bigcap_{i\in s}\varsigma(\varphi_i)
  \cap
  \bigcap_{i\in I\setminus s}\varsigma(\neg\varphi_i).
\]
It is therefore open in $\tau_{\Phi^{\pm}}$.  Every open generated by $\Phi^{\pm}$
is saturated under equality of signatures, so it is a union of these fibers.
\end{proof}

\begin{example}[Two positive observations]
Take $\Phi=(X,Y)$ on $\Omega_{X,Y}$.  Then
\[
 \tau_\Phi=
 \{\varnothing,\{B\},\{A,B\},\{B,C\},
   \{A,B,C\},\Omega_{X,Y}\}.
\]
Adding $\neg X$ and $\neg Y$ makes each of $A,B,C,D$ open and gives the
discrete topology.
\end{example}

\section{Distinguishability and separation}

Three questions should not be conflated:
\begin{enumerate}
\item Formulas $\mathcal L_1,\mathcal L_2$ are semantically different when
      $\varsigma(\mathcal L_1)\neq\varsigma(\mathcal L_2)$.
\item Valuations $v,w$ are observationally distinguishable by $\Phi$ when
      $o_\Phi(v)\neq o_\Phi(w)$.
\item Points are topologically separated according to the neighborhoods
      available in $\tau_\Phi$.
\end{enumerate}
Overlap between two truth regions does not by itself establish topological
indistinguishability.

\begin{proposition}[Semantic XOR criterion]
For formulas $\mathcal L_1,\mathcal L_2$,
\[
 \varsigma(\mathcal L_1)\neq\varsigma(\mathcal L_2)
 \quad\Longleftrightarrow\quad
 \varsigma(\mathcal L_1\Updownarrow\mathcal L_2)\neq\varnothing.
\]
At a particular valuation $v$, their truth values are distinct exactly when
$v\in\varsigma(\mathcal L_1\Updownarrow\mathcal L_2)$.
\end{proposition}

\begin{proof}
The XOR truth region is the symmetric difference
\[
 \varsigma(\mathcal L_1)\mathbin{\triangle}
 \varsigma(\mathcal L_2).
\]
It is nonempty exactly when the two truth regions are unequal.
\end{proof}

Observing only $\mathcal L_1\wedge\mathcal L_2$ gives one positive bit and
groups together all three ways in which the conjunction can be false.  Replacing
that observation by the pair $\Phi=(\mathcal L_1,\mathcal L_2)$ refines those
states according to the realized truth vectors $(1,1),(1,0),(0,1),(0,0)$.
Adding both negations makes the nonempty truth-pattern cells clopen.  This is
the precise observational sense in which the conjuncts become distinguishable.

\begin{definition}[Topological indistinguishability]
Points $v,w\in\Omega_P$ are topologically indistinguishable when they belong to
exactly the same open sets of $\tau_\Phi$.
\end{definition}

\begin{theorem}[Indistinguishability theorem]
For $v,w\in\Omega_P$, the following are equivalent:
\begin{enumerate}[label=(\roman*)]
\item $v$ and $w$ are topologically indistinguishable in $\tau_\Phi$;
\item $v$ and $w$ have the same subbasic open memberships;
\item $o_\Phi(v)=o_\Phi(w)$;
\item $v\sim_\Phi w$.
\end{enumerate}
\end{theorem}

\begin{proof}
Equality of all open memberships implies equality of subbasic memberships.
Subbasic memberships are exactly the coordinates of $o_\Phi$.  Conversely,
equality of signatures gives equality of membership in every finite
intersection of subbasic opens and hence in every union of such intersections.
\end{proof}

\begin{corollary}[$T_0$ quotient]
The space $(\Omega_P,\tau_\Phi)$ is $T_0$ exactly when $o_\Phi$ is injective.  Its
Kolmogorov quotient is naturally identified with the set $o_\Phi(\Omega_P)$ of
realized signatures, equipped with the subspace upper-set topology inherited
from $\mathbb S^I$.
\end{corollary}

The next separation level requires distinguishability in both directions.

\begin{proposition}[$T_1$ and complete finite separation]
The observation space is $T_1$ exactly when, for every distinct $v,w$,
\[
  o_\Phi(v)\nsubseteq o_\Phi(w)
  \quad\text{and}\quad
  o_\Phi(w)\nsubseteq o_\Phi(v).
\]
Equivalently, $o_\Phi$ is injective and its image is an antichain.  A
finite $T_1$ space is discrete; therefore finite Hausdorff observation spaces
are discrete as well.
\end{proposition}

\begin{proof}
There is an open containing $v$ but not $w$ exactly when some observation is
true at $v$ and false at $w$, which is equivalent to
$o_\Phi(v)\nsubseteq o_\Phi(w)$.  The $T_1$ condition requires this in both
directions.  In a finite $T_1$ space, each singleton is the finite intersection
of open complements of the other singletons, so every singleton is open.
\end{proof}

\begin{example}[Sierpi\'nski observation space]
Let $P=\{X\}$ and $\Phi=(X)$.  With valuations $0,1$, the topology is
\[
  \tau_\Phi=\{\varnothing,\{1\},\{0,1\}\}.
\]
It is $T_0$ but not $T_1$.  The positive observation $X$ distinguishes $1$
from $0$, but no positive observation in $\Phi$ distinguishes $0$ from $1$.
This is the canonical two-point example of asymmetric observation.
\end{example}

\section{Asymmetric separation and the specialization preorder}

We use the specialization convention
\[
 v\preceq_\Phi w
 \quad\Longleftrightarrow\quad
 \text{every open neighborhood of $v$ contains $w$}.
\]
Some authors use the reverse convention; the direction is fixed here by the
display above.

Because $\Omega_P$ is finite, every point has a smallest open neighborhood
\[
 N_\Phi(v):=\bigcap\{U\in\tau_\Phi:v\in U\}.
\]

\begin{theorem}[Positive-information order]
For all $v,w\in\Omega_P$,
\[
 v\preceq_\Phi w
 \quad\Longleftrightarrow\quad
 o_\Phi(v)\subseteq o_\Phi(w).
\]
Moreover,
\[
 N_\Phi(v)
 =\bigcap_{i\in o_\Phi(v)}\varsigma(\varphi_i)
 =\{w:o_\Phi(v)\subseteq o_\Phi(w)\}.
\]
\end{theorem}

\begin{proof}
If every open containing $v$ contains $w$, this holds for every subbasic open,
so every observation true at $v$ is true at $w$.  Conversely, if
$o_\Phi(v)\subseteq o_\Phi(w)$, then every basic open containing $v$ contains $w$,
and therefore every open containing $v$ contains $w$.  The neighborhood formula
follows.
\end{proof}

Thus $w$ preserves all positive information observable at $v$ and may satisfy
additional observations.  Failure of symmetry is not a defect: it records that
positive evidence can separate one direction without separating the reverse.
After quotienting by $\sim_\Phi$, the preorder becomes the inclusion order on
realized signatures.

\section{XOR difference and logical distance}

Coordinatewise XOR naturally produces a difference vector.  Let
$\mathbf a=(a_i)_{i\in I}$ be positive real weights.  As noted in Section~3,
these formulas use complete semantic profiles, even when the topology is read as
an interface for affirmative certificates.

\begin{definition}[Observable XOR difference]
Define
\[
 \Delta_\Phi(v,w)
 :=o_\Phi(v)\Updownarrow o_\Phi(w)
 :=\bigl(b_i(v)\Updownarrow b_i(w)\bigr)_{i\in I}.
\]
This vector identifies exactly which observations disagree.
\end{definition}

\begin{definition}[Symmetric logical distance]
Define the weighted Hamming distance between signatures by
\[
 d_\Phi(v,w)
 :=\sum_{i\in I} a_i\lvert b_i(v)-b_i(w)\rvert.
\]
This is a standard weighted form of Hamming distance \cite{Hamming1950}.
\end{definition}

\begin{proposition}
The function $d_\Phi$ is a pseudometric on $\Omega_P$, and
\[
 d_\Phi(v,w)=0
 \quad\Longleftrightarrow\quad
 v\sim_\Phi w.
\]
Consequently, $d_\Phi$ induces a metric on the underlying set
$\Omega_P/{\sim_\Phi}$.  If $\tau_{d_\Phi}$ denotes the pseudometric topology on
$\Omega_P$, then
\[
 \tau_{d_\Phi}=\tau_{\Phi^{\pm}}.
\]
The induced metric topology on the finite quotient set is discrete and need not
equal its Kolmogorov quotient topology.
\end{proposition}

\begin{proof}
Each coordinate contributes the weighted discrete metric on $\{0,1\}$; their
sum is a pseudometric.  It vanishes exactly when every observation coordinate
agrees.  If $I\neq\varnothing$, every sufficiently small ball at $v$ is exactly the signature
fiber $[v]_\Phi$, while every ball is a union of such fibers.  Hence the
pseudometric topology is the partition topology of Proposition
\ref{prop:two-sided-observation}.  For $I=\varnothing$ there is one fiber and both
topologies are indiscrete.  A genuine metric on a finite set has discrete
topology, giving the final qualification.
\end{proof}

For example, the Sierpi\'nski space of Example 4.6 is already $T_0$, so its
Kolmogorov quotient has the same nondiscrete topology.  Its Hamming metric is
the ordinary discrete two-point metric.  Thus ``metric on the quotient set''
does not mean ``metric inducing the quotient topology.''

Symmetric distance measures disagreement but cannot by itself encode the
direction of the specialization preorder.  For that purpose define the cost of
losing observations.

\begin{definition}[Directed logical distance]
Define
\[
 d_\Phi^+(v,w)
 :=\sum_{i\in I} a_i\,
   \ind[b_i(v)=1\text{ and }b_i(w)=0].
\]
Equivalently, this is weighted set difference from the support of $v$ to the
support of $w$.  In XOR notation the direction is supplied by the source mask:
\[
 d_\Phi^+(v,w)=\sum_{i\in I} a_i b_i(v)
   \bigl(b_i(v)\Updownarrow b_i(w)\bigr).
\]
This finite weighted sum is an elementary directed set-difference construction.
Pavlovic uses the cardinality $\lvert y\setminus x\rvert$, with the opposite
argument orientation from our loss convention \cite[Sec.~2]{Pavlovic2012},
while de Brecht uses a supremum of decaying weights on a countable powerset
\cite[Sec.~4]{deBrecht2013}.
\end{definition}

\begin{theorem}[Quasi-pseudometric representation]
The function $d_\Phi^+$ is a quasi-pseudometric: it is nonnegative,
$d_\Phi^+(v,v)=0$, and
\[
 d_\Phi^+(v,u)\leq d_\Phi^+(v,w)+d_\Phi^+(w,u).
\]
It also satisfies
\[
 d_\Phi^+(v,w)=0
 \quad\Longleftrightarrow\quad
 v\preceq_\Phi w,
\]
and
\[
 d_\Phi(v,w)=d_\Phi^+(v,w)+d_\Phi^+(w,v).
\]
The topology generated by the forward balls, for $r>0$,
\[
 B^+(v,r):=\{w:d_\Phi^+(v,w)<r\}
\]
is exactly $\tau_\Phi$.
\end{theorem}

\begin{proof}
For each coordinate, an observation lost from $v$ to $u$ must be lost either
from $v$ to $w$ or from $w$ to $u$.  Summing the resulting coordinatewise
inequality proves the triangle inequality.  Vanishing means that no observation
true at $v$ is false at $w$, which is exactly
$o_\Phi(v)\subseteq o_\Phi(w)$.  The symmetrization identity follows because every
disagreement is a loss in exactly one direction.

Every forward ball is upward closed in $\preceq_\Phi$: if $w$ belongs to the ball
and $w\preceq_\Phi u$, then $u$ loses no more observations from $v$ than $w$
does.  Hence every forward ball is $\tau_\Phi$-open.  If $I\neq\varnothing$ and
$0<r\leq\min_{i\in I} a_i$, then
\[
 B^+(v,r)=N_\Phi(v).
\]
The smallest neighborhoods generate $\tau_\Phi$, so the forward-ball topology is
exactly $\tau_\Phi$.  If $I=\varnothing$, the directed distance is identically zero and
every positive-radius forward ball is $\Omega_P$, again giving the correct
indiscrete topology.
\end{proof}

The zero-distance relation of $d_\Phi^+$ recovers the preorder and its forward
balls recover the topology.  Its positive numerical values contain additional
presentation data: changing weights, or adding a topologically redundant
coordinate, can change distances without changing $\preceq_\Phi$ or
$\tau_\Phi$.  This is compatible with the general role of nonsymmetric distance
in generalized metric theory \cite{Lawvere1973}.

\section{Interior, closure, and observational boundary}

Topological operators describe what can be verified or ruled out using the
chosen observations.

\begin{definition}
For an event $E\subseteq\Omega_P$, let
\[
 \Int_\Phi(E):=\Int_{\tau_\Phi}(E),\qquad
 \Cl_\Phi(E):=\Cl_{\tau_\Phi}(E),\qquad
 \bd_\Phi E:=\Cl_\Phi(E)\setminus\Int_\Phi(E).
\]
\end{definition}

The interior $\Int_\Phi(E)$ is the largest observable open event contained in
$E$.  If the actual valuation lies in this interior, a finite conjunction of
available positive observations certifies membership in $E$.  The closure $\Cl_\Phi(E)$ consists of
valuations at which the available observations cannot exclude contact with
$E$.  The boundary $\bd_\Phi E$ is therefore an observational ambiguity region.
Equivalently,
\[
 \bd_\Phi E
 =\Cl_\Phi(E)\cap\Cl_\Phi(\Omega_P\setminus E).
\]
It contains precisely those points whose available neighborhoods cannot settle
membership in $E$ one way or the other.

\begin{proposition}[Finite neighborhood formulas]
For every $E\subseteq\Omega_P$,
\begin{align*}
 \Int_\Phi(E)&=\{v:N_\Phi(v)\subseteq E\},\\
 \Cl_\Phi(E)&=\{v:N_\Phi(v)\cap E\neq\varnothing\}.
\end{align*}
Consequently,
\[
 v\in\bd_\Phi E
 \quad\Longleftrightarrow\quad
 N_\Phi(v)\cap E\neq\varnothing
 \text{ and }
 N_\Phi(v)\cap(\Omega_P\setminus E)\neq\varnothing.
\]
\end{proposition}

\begin{proof}
The point $v$ is interior to $E$ exactly when it has an open neighborhood
contained in $E$.  Since $N_\Phi(v)$ is its smallest neighborhood, this is
equivalent to $N_\Phi(v)\subseteq E$.  A point belongs to the closure of $E$
exactly when each of its neighborhoods meets $E$, which in a finite space is
equivalent to $N_\Phi(v)\cap E\neq\varnothing$.
\end{proof}

\begin{remark}[Boundary versus XOR]
The symbol $\bd_\Phi$ denotes the topological boundary operator above.  It is
not the Boolean symmetric-difference operation that removes the overlap of two
formulas.  That Boolean operation is the exclusive remainder
\[
 \varsigma(\mathcal L_1\vee\mathcal L_2)
 \setminus\varsigma(\mathcal L_1\wedge\mathcal L_2)
 =\varsigma(\mathcal L_1\Updownarrow\mathcal L_2).
\]
The two constructions can differ even in the simplest cases: in a discrete
topology every topological boundary is empty, while an XOR truth region
may be nonempty.
\end{remark}

\begin{example}[Boundary of implication under positive observations]
For $\Phi=(X,Y)$,
\[
 E=\varsigma(X\Rightarrow Y)=\{B,C,D\}.
\]
Using the topology in Example 3.6,
\[
 \Int_\Phi(E)=\{B,C\},\qquad
 \Cl_\Phi(E)=\Omega_{X,Y},\qquad
 \bd_\Phi E=\{A,D\}.
\]
Thus $X\Rightarrow Y$ is neither open nor closed in this observation topology.
The point $D$ satisfies the implication but is not positively certifiable as
such using only $X$ and $Y$, while $A$ lies observationally adjacent to
satisfying states despite falsifying the implication.
\end{example}

\begin{example}[XOR region under coarse and discrete observation]
Retain the partition labels $A=X\wedge\neg Y$,
$B=X\wedge Y$, $C=\neg X\wedge Y$, and
$D=\neg X\wedge\neg Y$.  The exclusive truth region is
\[
 E_\oplus:=\varsigma(X\Updownarrow Y)=\{A,C\}.
\]
First use only the positive observations $\Phi=(X,Y)$.  The smallest
neighborhoods are
\[
 N_\Phi(A)=\{A,B\},\quad N_\Phi(B)=\{B\},\quad
 N_\Phi(C)=\{B,C\},\quad N_\Phi(D)=\Omega_{X,Y}.
\]
Hence no nonempty $\tau_\Phi$-open set is contained in $E_\oplus$, while the
neighborhoods of $A$, $C$, and $D$ each meet it.  Thus
\[
 \Int_\Phi(E_\oplus)=\varnothing,\qquad
 \Cl_\Phi(E_\oplus)=\{A,C,D\},\qquad
 \bd_\Phi E_\oplus=\{A,C,D\}.
\]
The points $A$ and $C$ belong to the XOR region but are not positively
certifiable using only $X$ and $Y$, because their smallest neighborhoods also
contain $B$.  The point $D$ is not in the XOR region, but it remains on its
boundary because every positive-observation neighborhood of $D$ contains both
XOR and non-XOR valuations.

Now take the two-sided observation family
\[
 \Phi^{\pm}=(X,\neg X,Y,\neg Y).
\]
Its topology is discrete.  Therefore $E_\oplus$ is both open and closed, and
\[
 \Int_{\Phi^{\pm}}(E_\oplus)
 =\Cl_{\Phi^{\pm}}(E_\oplus)
 =\{A,C\},
 \qquad
 \bd_{\Phi^{\pm}}E_\oplus=\varnothing.
\]
The XOR truth set is unchanged by the topology; only its observable boundary
changes.  This is the distinction between Boolean symmetric difference and topological
ambiguity.  Figure~\ref{fig:xor-four-cell} displays the same calculation on
the four truth-pattern cells.
\end{example}

\begin{figure}[!ht]
\centering
\renewcommand{\arraystretch}{1.35}
\setlength{\tabcolsep}{13pt}
\begin{tabular}{c|c|c}
 & $X$ & $\neg X$ \\\hline
$Y$ &
$\begin{gathered}B=X\wedge Y\\\text{non-XOR}\end{gathered}$ &
$\begin{gathered}C=\neg X\wedge Y\\\mathbf{XOR}\end{gathered}$ \\\hline
$\neg Y$ &
$\begin{gathered}A=X\wedge\neg Y\\\mathbf{XOR}\end{gathered}$ &
$\begin{gathered}D=\neg X\wedge\neg Y\\\text{non-XOR}\end{gathered}$
\end{tabular}
\caption{The four truth-pattern cells for $X$ and $Y$.  The bold cells form
$E_\oplus=\varsigma(X\Updownarrow Y)=\{A,C\}$.  Under positive observation
$\Phi=(X,Y)$, the observable boundary is $\bd_\Phi E_\oplus=\{A,C,D\}$;
under $\Phi^{\pm}$ it is empty.}
\label{fig:xor-four-cell}
\end{figure}

\subsection{Open truth regions and Heyting implication}

The opens of any topology form a Heyting algebra.  For
$U,V\in\tau_\Phi$, define
\[
 U\Rightarrow_H V
 :=\Int_\Phi\bigl((\Omega_P\setminus U)\cup V\bigr),
 \qquad
 \neg_H U:=\Int_\Phi(\Omega_P\setminus U).
\]
Then
\[
 U\cap(U\Rightarrow_H V)\subseteq V
\]
is topological Modus Ponens, and
\[
 (U\Rightarrow_H V)\cap\neg_H V\subseteq\neg_H U
\]
is topological Modus Tollens.  In a discrete topology the interiors do nothing,
and these reduce to their classical truth-set forms.  In a nondiscrete
observation topology, the interior records the loss imposed by requiring the
result to remain observable.

Indeed, for every open $W$,
\[
 W\subseteq U\Rightarrow_H V
 \quad\Longleftrightarrow\quad
 W\cap U\subseteq V,
\]
which is the defining Heyting adjunction.  The two displayed inference
inequalities follow from it.  These open-set operations belong to standard
topological semantics \cite{McKinseyTarski1944,Vickers2007}; they should be
distinguished from the classical truth-region operations used earlier.  A full intuitionistic
interpretation would assign open sets to atoms and evaluate compound formulas
recursively in this Heyting algebra.

For example, in the Sierpi\'nski space of Example~4.6, interpret $X$ by the
open $\{1\}$.  Then $\neg_H\{1\}=\varnothing$ and
$\neg_H\neg_H\{1\}=\Omega_X$, whereas the interior of the classical truth
region $\varsigma(\neg\neg X)=\{1\}$ is still $\{1\}$.  Recursive Heyting
evaluation therefore differs from taking one final interior of a classical
truth region.

\section{Information refinement and logical update}

The phrase ``acquiring information'' can describe two mathematically different
operations.

\subsection{Refining the observations}

Let $\Phi=(\varphi_i)_{i\in I}$ and $\Psi=(\psi_j)_{j\in J}$, where
$I\subseteq J$ and $\psi_i=\varphi_i$ for every $i\in I$.  Thus $\Psi$ extends
$\Phi$ while retaining the identities of its old coordinates.

\begin{theorem}[Refinement theorem]
For such an indexed extension $\Phi\hookrightarrow\Psi$:
\begin{enumerate}[label=(\roman*)]
\item $\tau_\Phi\subseteq\tau_\Psi$;
\item $v\sim_\Psi w$ implies $v\sim_\Phi w$;
\item $\preceq_\Psi\,\subseteq\,\preceq_\Phi$ as relations;
\item if $\Psi$ retains the weights $a_i$ on $i\in I$, then
      $d_\Psi^+(v,w)\geq d_\Phi^+(v,w)$ and $d_\Psi(v,w)\geq d_\Phi(v,w)$.
\end{enumerate}
\end{theorem}

\begin{proof}
The subbasis for $\tau_\Phi$ is contained in the subbasis for $\tau_\Psi$, proving
(i).  Equality of the longer signatures implies equality after deleting the
new coordinates, proving (ii).  Inclusion of $\Psi$-signatures implies inclusion
after deleting those coordinates, proving (iii).  The added coordinates make
nonnegative contributions to both distances, proving (iv).
\end{proof}

Adding an observable can split an indistinguishability block, remove a
one-way comparison, and increase measured separation.  A proposed observation
$\psi$ is topologically redundant relative to $\Phi$ exactly when
\[
 \varsigma(\psi)\in\tau_\Phi,
\]
because adding an already open truth region does not refine the topology.
Topological redundancy need not mean numerical redundancy: an added coordinate
can add a further term to $d_\Phi$ and $d_\Phi^+$ while leaving the generated
topology unchanged.  For example, adding $X\wedge Y$ to $\Phi=(X,Y)$ is
topologically redundant because $\varsigma(X\wedge Y)$ is already open.  With
unit weights it nevertheless changes $d^+(B,A)$ from $1$ to $2$.

\subsection{Updating the possible valuations}

Let $S\subseteq\Omega_P$ be the current possibility region.  Learning a
premise $\gamma$ produces
\[
 S\longmapsto S\cap\varsigma(\gamma).
\]
For premises $\Gamma=\{\gamma_1,\ldots,\gamma_r\}$, define
\[
 S[\Gamma]:=S\cap\bigcap_{j=1}^r\varsigma(\gamma_j),
 \qquad
 S_\Gamma:=\Omega_P[\Gamma].
\]
For an arbitrary current carrier $S$, the updated observation space is
\[
 \mathbf X^\True_{S,\Gamma}
 :=\bigl(S[\Gamma],\tau_\Phi\restr_{S[\Gamma]}\bigr),
 \qquad
 \tau_\Phi\restr_S:=\{U\cap S:U\in\tau_\Phi\}.
\]
When $S=\Omega_P$, write
$\mathbf X^\True_\Gamma:=\mathbf X^\True_{\Omega_P,\Gamma}$.
Sequential updates are associative, commutative, and idempotent because set
intersection has those properties.  Inconsistent premises produce the empty
space.  Smallest neighborhoods restrict accordingly:
\[
 N_{\Phi\mid S[\Gamma]}(v):=S[\Gamma]\cap N_\Phi(v)
 \qquad(v\in S[\Gamma]).
\]

Observation refinement and premise update need not have the same effect.
Refinement changes which distinctions are expressible; update changes which
worlds remain possible.

\section{Inference by truth-region restriction}

Classical inference is a relationship between the common truth region of the
premises and the truth region of the conclusion.

\begin{theorem}[Semantic consequence]
For a finite set of premises $\Gamma$ and a conclusion $\psi$,
\[
 \Gamma\models\psi
 \quad\Longleftrightarrow\quad
 S_\Gamma\subseteq\varsigma(\psi).
\]
When this holds, the inclusion
\[
 \mathbf X^\True_\Gamma
 \hookrightarrow
 \bigl(\varsigma(\psi),
       \tau_\Phi\restr_{\varsigma(\psi)}\bigr)
\]
is continuous.
\end{theorem}

\begin{proof}
The semantic statement says exactly that every valuation satisfying all
premises satisfies $\psi$.  The inclusion between two subspaces of the same
ambient space is continuous.
\end{proof}

The topology contributes a structured representation of this consequence: the
common premise region carries an induced observation topology, and the valid
inclusion is a continuous map into the conclusion subspace.  It does not alter
which consequences are classically valid.  In particular, premise conjunction
uses intersection of truth regions (or the pullback of their induced
subspaces), not intersection of two families of open sets.  The latter is a
different operation: a meet in the lattice of topologies on one fixed carrier.

\subsection{A remark on minimal truth spaces}

For a truth region $S$, one may form the indiscrete space
$\mathcal T^\True_S=(S,\{\varnothing,S\})$.  Intersecting carriers and then
assigning this indiscrete topology transports the ordinary meet
$R\cap S$ of the power set into new notation; it supplies no additional
topological inference principle.  Moreover, $\mathcal T^\True_S$ generally
differs from the induced observation subspace
$(S,\tau_\Phi\restr_S)$, and its inclusion into the ambient observation space
need not be continuous.  It therefore cannot replace the induced subspaces in
the pullback formulation below.  The two spaces coincide exactly when the
induced topology on $S$ is indiscrete, as for a singleton or a region whose
points all have the same observation signature.

\subsection{Modus Ponens and Modus Tollens}

Using $A,B,C,D$,
\begin{align*}
 \varsigma(X\Rightarrow Y)\cap\varsigma(X)
   &=\{B\}
    =\varsigma(X\wedge Y)
    \subseteq\varsigma(Y),\\
 \varsigma(X\Rightarrow Y)\cap\varsigma(\neg Y)
   &=\{D\}
    =\varsigma(\neg X\wedge\neg Y)
    \subseteq\varsigma(\neg X).
\end{align*}
\subsection{Hypothetical Syllogism}

For three variables use the partition labels
\[
\begin{array}{llll}
A=X\wedge\neg Y\wedge\neg Z,
&B=\neg X\wedge Y\wedge\neg Z,
&C=\neg X\wedge\neg Y\wedge Z,
&D=X\wedge Y\wedge\neg Z,\\
E=X\wedge\neg Y\wedge Z,
&F=\neg X\wedge Y\wedge Z,
&G=X\wedge Y\wedge Z,
&H=\neg X\wedge\neg Y\wedge\neg Z.
\end{array}
\]
Then
\begin{align*}
 \varsigma(X\Rightarrow Y)&=\{B,C,D,F,G,H\},\\
 \varsigma(Y\Rightarrow Z)&=\{A,C,E,F,G,H\},\\
 \varsigma(X\Rightarrow Z)&=\{B,C,E,F,G,H\}.
\end{align*}
Hence
\[
 \varsigma(X\Rightarrow Y)\cap\varsigma(Y\Rightarrow Z)
 =\{C,F,G,H\}
 \subsetneq\varsigma(X\Rightarrow Z).
\]
Here the inclusion is proper because the premises contain more information than
the conclusion.  Validity in general requires inclusion, which may also be an
equality.
Figure~\ref{fig:hypothetical-syllogism} makes the truth-region calculation
explicit: the boxed cells satisfy both premises and every one satisfies the
conclusion.

\begin{figure}[!ht]
\centering
\renewcommand{\arraystretch}{1.22}
\setlength{\tabcolsep}{10pt}
\begin{tabular}{c|c|c|c|c}
\text{cell} & $X\Rightarrow Y$ & $Y\Rightarrow Z$ &
\text{both premises} & $X\Rightarrow Z$ \\\hline
$A$ & $\False$ & $\True$  & $\False$ & $\False$ \\
$B$ & $\True$  & $\False$ & $\False$ & $\True$  \\
$C$ & $\True$  & $\True$  & $\boxed{\True}$ & $\True$  \\
$D$ & $\True$  & $\False$ & $\False$ & $\False$ \\
$E$ & $\False$ & $\True$  & $\False$ & $\True$  \\
$F$ & $\True$  & $\True$  & $\boxed{\True}$ & $\True$  \\
$G$ & $\True$  & $\True$  & $\boxed{\True}$ & $\True$  \\
$H$ & $\True$  & $\True$  & $\boxed{\True}$ & $\True$
\end{tabular}
\caption{Hypothetical syllogism on the eight three-variable truth-pattern
cells.  The boxed entries are
$\varsigma(X\Rightarrow Y)\cap\varsigma(Y\Rightarrow Z)=\{C,F,G,H\}$;
they all lie in $\varsigma(X\Rightarrow Z)$.  The operation is an intersection
of truth regions, equivalently the pullback of their truth subspaces, not an
intersection of the families of open sets.}
\label{fig:hypothetical-syllogism}
\end{figure}

\FloatBarrier
\subsection{Pullback formulation}

For a single formula $\mathcal L$, write
\[
 \mathbf X^\True_{\mathcal L}
 :=\bigl(\varsigma(\mathcal L),
          \tau_\Phi\restr_{\varsigma(\mathcal L)}\bigr).
\]
The inclusions of two truth subspaces into $(\Omega_P,\tau_\Phi)$ have pullback
\[
 \mathbf X^\True_{\mathcal L_1}
 \times_{(\Omega_P,\tau_\Phi)}
 \mathbf X^\True_{\mathcal L_2}
 \cong
 \mathbf X^\True_{\mathcal L_1\wedge\mathcal L_2}.
\]
For Modus Ponens,
\[
 \mathbf X^\True_{X\Rightarrow Y}
 \times_{(\Omega_{X,Y},\tau_\Phi)}
 \mathbf X^\True_X
 \cong
 \mathbf X^\True_{X\wedge Y}
 \hookrightarrow
 \mathbf X^\True_Y.
\]
This is the standard topological construction closest to the idea of forming
a relative space from several premises.  Explicitly, its points are pairs
$(v,v)$ with $v$ in both truth regions, and the diagonal identification with the
intersection is a homeomorphism.  It also has the usual universal property:
any two continuous maps with equal composites into the ambient observation
space factor uniquely through this intersection subspace.  Indeed, if
$f:Z\to\varsigma(\mathcal L_1)$ and
$g:Z\to\varsigma(\mathcal L_2)$ have equal ambient composites, define $h(z)$
to be their common value.  For an ambient open $U$,
\[
 h^{-1}\!\left(U\cap\varsigma(\mathcal L_1)\cap
 \varsigma(\mathcal L_2)\right)
 =f^{-1}\!\left(U\cap\varsigma(\mathcal L_1)\right),
\]
which is open; hence the unique set-theoretic mediator is continuous.

\section{Continuous maps and preservation of information}

Let $\Phi=(\varphi_i)_{i\in I}$ and $\Psi=(\psi_j)_{j\in J}$ generate the
observation spaces $(\Omega_P,\tau_\Phi)$ and $(\Omega_Q,\tau_\Psi)$.  A map
\[
 f:\Omega_P\longrightarrow\Omega_Q
\]
is continuous when the inverse image of every $\Psi$-observable open event is
$\Phi$-observable.  It suffices to check the subbasic truth regions:
\[
 f^{-1}\bigl(\varsigma(\psi_j)\bigr)\in\tau_\Phi
 \qquad(j\in J).
\]

\begin{proposition}[Monotonicity of continuous maps]
If $f$ is continuous and $v\preceq_\Phi w$, then
\[
 f(v)\preceq_\Psi f(w).
\]
\end{proposition}

\begin{proof}
Let $U\in\tau_\Psi$ contain $f(v)$.  Then $f^{-1}(U)$ is a $\Phi$-open set
containing $v$, so it contains $w$.  Hence $U$ contains $f(w)$.
\end{proof}

Continuity therefore preserves one-way observational dependence.  This gives
a rigorous starting point for continuous valuation maps.  Conversely, because
these finite spaces have their upper-set topologies, every monotone valuation
map is continuous.  Continuity may collapse distinguishable points and does not
by itself imply preservation of either weighted distance.

\section{Scope, limitations, and deferred directions}

The results above form a finite expository theory of observation,
distinguishability, distance, refinement, and inference.  They do not establish
new results in homology, sheaf theory, cognition, or quantum theory.  The most
direct extensions are questions about the design and use of observations in
reasoning.

\subsection{Logical complexes and homology}

A standard relation-nerve construction can be applied to a family of truth
regions.  For an indexed observation family
$\Phi=(\varphi_i)_{i\in I}$, define a simplicial complex
\[
 K_\Phi:=\left\{J\subseteq I:
       \bigcap_{i\in J}\varsigma(\varphi_i)\neq\varnothing
       \right\}.
\]
Its vertices are the indices of individually satisfiable observations, and its
simplices are jointly satisfiable indexed subfamilies.  This is the
observation-side Dowker complex
of the satisfaction relation \cite{Dowker1952}.  It can record higher-order
compatibility not visible from pairwise intersections, but it is not an
invariant of $\tau_\Phi$.  For example,
$\Phi=(X,Y,\neg(X\wedge Y))$ gives the boundary of a triangle.  Its observation
regions are $\{A,B\}$, $\{B,C\}$, and $\{A,C,D\}$: their pairwise intersections
are respectively $\{B\}$, $\{A\}$, and $\{C\}$, while their triple intersection
is empty.  Thus the nerve is connected, whereas the observation space has
connected components $\{B\}$ and $\{A,C,D\}$.  Adjoining a
tautology leaves the generated topology and both signature distances unchanged,
but makes the nerve a cone and kills its reduced homology.  Any homological use
must therefore state the presentation and the intended invariance separately:
the homology of $K_\Phi$ need not agree with the singular homology of the
observation space and is not determined by $\tau_\Phi$.

\subsection{Reasoning-oriented directions}

Concrete next questions include finding a smallest or least-cost observation
family that separates designated valuations, certifies a target conclusion, or
distinguishes competing premise sets; measuring the complexity of those design
problems; and studying sequences in which observation refinement and premise
restriction alternate.  Another direction is to compare the classical
truth-region semantics with the recursively defined Heyting semantics of the
open-set algebra, or with other nonclassical logics whose evidence conditions
are explicit.

A contextual extension would additionally require a cover by reasoning
contexts, specified sets of allowed local assignments, and restrictions on
overlaps.  Nonexistence of any support-compatible global assignment is a
strong-contextuality-style obstruction.  General probabilistic contextuality
instead asks for a global distribution reproducing the local distributions, so
the existence of one compatible assignment does not decide that broader
question.  The unrestricted assignment sheaf itself already glues compatible
assignments \cite[Sec.~6]{AbramskyBrandenburger2011}.

Other finite-space properties, such as connectedness, can be read from the
specialization preorder.  Form the undirected graph with an edge between
distinct points $v,w$ whenever $v\preceq_\Phi w$ or $w\preceq_\Phi v$, including
distinct equivalent points.  Its graph components are exactly the topological
connected components, without a $T_0$ assumption: each graph component is both
an upper and a lower set and hence clopen, while a clopen separation permits no
comparison across it \cite[Prop.~5]{Stong1966}.  This is a useful descriptive
corollary, but a full catalogue of standard finite-space invariants would obscure the paper's
observation-versus-premise focus.

\section{Conclusion}

A propositional logical topology can be nontrivial when it records a selected
family of observations rather than every definable truth set.  Observation
signatures simultaneously determine a partition into indistinguishable states,
a finite topology, and its specialization preorder.  Coordinatewise XOR records
which observations differ; weighted Hamming distance measures symmetric
disagreement; and the directed distance $d_\Phi^+$ measures loss of positive
information while generating the selected observation topology.  Adding tests
refines distinguishability, while adding premises restricts the possibility
space.  Inference is the inclusion of that restricted region in the truth
region of a conclusion.

The framework gives a common notation for partitioning, distinguishability,
asymmetric separation, logical distance, and refinement.  Its individual
ingredients are classical.  Its value here is expository: it makes explicit
which structures are determined by the observations, which require extra
weights or presentation choices, and which logical claims remain ordinary
model-set inclusions.  A further research contribution would require an
additional algorithm, complexity result, invariant, or validated application.

\begin{thebibliography}{99}
\small

\bibitem{AbramskyBrandenburger2011}
S. Abramsky and A. Brandenburger,
``The sheaf-theoretic structure of non-locality and contextuality,''
\emph{New Journal of Physics} 13 (2011), 113036.
\href{https://doi.org/10.1088/1367-2630/13/11/113036}
{doi:10.1088/1367-2630/13/11/113036}.

\bibitem{Alexandroff1937}
P. Alexandroff,
``Diskrete R\"aume,''
\emph{Recueil Math\'ematique, Nouvelle S\'erie} 2(44), no. 3 (1937),
501--519.

\bibitem{BonsangueJacobsKok1995}
M. M. Bonsangue, B. P. F. Jacobs, and J. N. Kok,
``Duality beyond sober spaces: Topological spaces and observation frames,''
\emph{Theoretical Computer Science} 151, no. 1 (1995), 79--124.
\href{https://doi.org/10.1016/0304-3975(95)00048-2}
{doi:10.1016/0304-3975(95)00048-2}.

\bibitem{deBrecht2013}
M. de Brecht,
``Quasi-Polish spaces,''
\emph{Annals of Pure and Applied Logic} 164, no. 3 (2013), 356--381.
\href{https://doi.org/10.1016/j.apal.2012.11.001}
{doi:10.1016/j.apal.2012.11.001}.

\bibitem{Dowker1952}
C. H. Dowker,
``Homology groups of relations,''
\emph{Annals of Mathematics}, second series, 56, no. 1 (1952), 84--95.
\href{https://doi.org/10.2307/1969768}{doi:10.2307/1969768}.

\bibitem{GanterWille1999}
B. Ganter and R. Wille,
\emph{Formal Concept Analysis: Mathematical Foundations},
Springer, Berlin, 1999.
\href{https://doi.org/10.1007/978-3-642-59830-2}
{doi:10.1007/978-3-642-59830-2}.

\bibitem{Hamming1950}
R. W. Hamming,
``Error detecting and error correcting codes,''
\emph{Bell System Technical Journal} 29, no. 2 (1950), 147--160.
\href{https://doi.org/10.1002/j.1538-7305.1950.tb00463.x}
{doi:10.1002/j.1538-7305.1950.tb00463.x}.

\bibitem{Lawvere1973}
F. W. Lawvere,
``Metric spaces, generalized logic, and closed categories,''
\emph{Rendiconti del Seminario Matematico e Fisico di Milano} 43 (1973),
135--166.
\href{https://doi.org/10.1007/BF02924844}{doi:10.1007/BF02924844}.

\bibitem{McKinseyTarski1944}
J. C. C. McKinsey and A. Tarski,
``The algebra of topology,''
\emph{Annals of Mathematics}, second series, 45, no. 1 (1944), 141--191.
\href{https://doi.org/10.2307/1969080}{doi:10.2307/1969080}.

\bibitem{Pawlak1982}
Z. Pawlak,
``Rough sets,''
\emph{International Journal of Computer and Information Sciences} 11 (1982),
341--356.
\href{https://doi.org/10.1007/BF01001956}{doi:10.1007/BF01001956}.

\bibitem{Pavlovic2012}
D. Pavlovic,
``Quantitative concept analysis,'' in \emph{Formal Concept Analysis},
Lecture Notes in Computer Science 7278 (2012), 260--277.
\href{https://doi.org/10.1007/978-3-642-29892-9_24}
{doi:10.1007/978-3-642-29892-9\_24}.

\bibitem{Stone1936}
M. H. Stone,
``The theory of representations for Boolean algebras,''
\emph{Transactions of the American Mathematical Society} 40, no. 1 (1936),
37--111.
\href{https://doi.org/10.1090/S0002-9947-1936-1501865-8}
{doi:10.1090/S0002-9947-1936-1501865-8}.

\bibitem{Stong1966}
R. E. Stong,
``Finite topological spaces,''
\emph{Transactions of the American Mathematical Society} 123, no. 2 (1966),
325--340.
\href{https://doi.org/10.1090/S0002-9947-1966-0195042-2}
{doi:10.1090/S0002-9947-1966-0195042-2}.

\bibitem{Vickers1989}
S. Vickers,
\emph{Topology via Logic}, Cambridge Tracts in Theoretical Computer Science 5,
Cambridge University Press, 1989.
\href{https://doi.org/10.1017/CBO9780511625900}
{doi:10.1017/CBO9780511625900}.

\bibitem{Vickers2007}
S. Vickers,
``Locales and toposes as spaces,'' in \emph{Handbook of Spatial Logics},
Springer, 2007, 429--496.
\href{https://sjvickers.github.io/LocTopSpaces.pdf}
{Author preprint}.

\end{thebibliography}

\end{document}
